% Preamble of the keynote: fonts, colours, frame layout, and the example shared by the slides.

% Korean text. NanumGothic is a static font; the variable Noto fonts of this laptop fail in xdvipdfmx.
\usepackage{kotex}
\setsansfont      {NanumGothic}[BoldFont = NanumGothic Bold]
\setsanshangulfont{NanumGothic}[BoldFont = NanumGothic Bold]

% Mathematics in a serif face, so that I, l and 1 stay apart
\usefonttheme{professionalfonts}
\usepackage{unicode-math}
\setmathfont{STIX Two Math}

% Figures
\usepackage{tikz}
\usetikzlibrary{arrows.meta, decorations.pathreplacing}

% Colours: orange is importance and selected rows, blue is latency. Everything else is grey.
\definecolor{ink}    {HTML}{22262B}
\definecolor{muted}  {HTML}{8A9099}
\definecolor{dim}    {HTML}{6F7781}
\definecolor{hair}   {HTML}{C9CED4}
\definecolor{off}    {HTML}{E8EAED}
\definecolor{imp}    {HTML}{D9531E}
\definecolor{impsoft}{HTML}{F5D3C2}
\definecolor{lat}    {HTML}{2B5C8A}
\definecolor{latsoft}{HTML}{D2DFEB}

% Text in one of the colours. The names avoid \lat and \minus, which unicode-math defines as symbols.
\newcommand{\inimp} [1]{{\color{imp}#1}}
\newcommand{\inlat} [1]{{\color{lat}#1}}
\newcommand{\insoft}[1]{{\color{muted}#1}}

% NanumGothic has no minus sign; the one of the math font is used in text.
\newcommand{\minussign}{\ensuremath{-}}

% Frame layout: the section name above the title, the frame number in the corner.
\setbeamertemplate{navigation symbols}{}
\setbeamersize{text margin left = 12mm, text margin right = 12mm}
\setbeamercolor{normal text}{fg = ink, bg = white}

\setbeamertemplate{frametitle}{%
  \vspace{7.5mm}%
  {\scriptsize\color{muted}\insertsectionhead}\par
  \vspace{0.8mm}%
  {\fontsize{13pt}{16pt}\selectfont\bfseries\color{ink}\insertframetitle}\par
  \vspace{1.5mm}%
}

\setbeamertemplate{footline}{%
  \hfill{\scriptsize\color{muted}\insertframenumber}\hspace{12mm}\vspace{5mm}%
}

% Every figure is drawn on one canvas of the text width, so that what the slides share stays in place.
\newenvironment{canvas}
  {\begin{tikzpicture}\useasboundingbox (0, 0) rectangle (13.6, 5.45);}
  {\end{tikzpicture}}

% Shared drawing styles
\tikzset{
  cell/.style      = {draw = white, line width = 0.8pt},
  cell on/.style   = {cell, fill = imp},
  cell off/.style  = {cell, fill = off},
  num on/.style    = {text = white, font = \small\bfseries},
  num off/.style   = {text = dim, font = \small},
  label/.style     = {font = \footnotesize},
  axis/.style      = {draw = ink, line width = 0.5pt},
  thin line/.style = {draw = hair, line width = 0.5pt},
  bracket/.style   = {draw = muted, line width = 0.5pt},
  brace/.style     = {draw = muted, line width = 0.5pt, decorate, decoration = {brace, amplitude = 3.5pt}},
  tip/.style       = {-{Stealth[length = 4pt, width = 3.5pt]}},
}

% The example of the slides: N = 16 rows, exactly R = 8 of them selected.
% Importance v_i of each row
\def\toyV{8,3,0,4,9,4,2,1,4,1,8,9,1,9,1,1}

% Masks of the example. Tiles with B = 1 is Top-R, and with B = 8 it is one chunk.
\def\maskNone {0,0,0,0,0,0,0,0,0,0,0,0,0,0,0,0}
\def\maskTop  {1,0,0,1,1,1,0,0,1,0,1,1,0,1,0,0}
\def\maskTwo  {1,1,0,0,1,1,0,0,0,0,1,1,1,1,0,0}
\def\maskFour {0,0,1,1,1,1,0,0,0,0,1,1,1,1,0,0}
\def\maskEight{0,0,0,0,1,1,1,1,1,1,1,1,0,0,0,0}
\def\maskOpt  {0,0,0,1,1,1,1,0,0,0,1,1,1,1,0,0}

% The chunks of the Top-R mask: start / length / importance I[C] / latency T[|C|] / latency of the chunks before it
\def\toyChunks{0/1/8/5/0, 3/3/17/7/5, 8/1/4/5/12, 10/2/17/6/17, 13/1/9/5/23}

% Where the N boxes of the example sit on the canvas: left edge, bottom edge, size of one box
\def\boxleft{2.5}
\def\boxbase{4.75}
\def\boxsize{0.53}

% \drawmask{mask}{x}{y}{cell width}{cell height} draws the N boxes of a mask from (x, y).
\newcommand{\drawmask}[5]{%
  \foreach \bit [count = \k from 0] in #1 {
    \ifnum\bit=1
      \path[cell on]  ({#2 + \k * #4}, {#3}) rectangle ++({#4}, {#5});
    \else
      \path[cell off] ({#2 + \k * #4}, {#3}) rectangle ++({#4}, {#5});
    \fi
  }%
}

% \drawvalues{mask} draws the N boxes of the example with the importance of each row written in its box.
% In the mask, 1 is a selected row and 2 is a selected row drawn pale.
\newcommand{\drawvalues}[1]{%
  \foreach \val [count = \k from 0] in \toyV {
    \pgfmathtruncatemacro\bit{{#1}[\k]}
    \ifcase\bit
      \path[cell off]             ({\boxleft + \k * \boxsize}, \boxbase) rectangle ++(\boxsize, \boxsize);
      \node[num off]              at ({\boxleft + (\k + 0.5) * \boxsize}, {\boxbase + 0.5 * \boxsize}) {\val};
    \or
      \path[cell on]              ({\boxleft + \k * \boxsize}, \boxbase) rectangle ++(\boxsize, \boxsize);
      \node[num on]               at ({\boxleft + (\k + 0.5) * \boxsize}, {\boxbase + 0.5 * \boxsize}) {\val};
    \or
      \path[cell, fill = impsoft] ({\boxleft + \k * \boxsize}, \boxbase) rectangle ++(\boxsize, \boxsize);
      \node[num on, text = imp]   at ({\boxleft + (\k + 0.5) * \boxsize}, {\boxbase + 0.5 * \boxsize}) {\val};
    \fi
  }%
}

% \drawbrackets marks the chunks of the Top-R mask below the boxes.
\newcommand{\drawbrackets}{%
  \foreach \start/\len/\importance/\latency/\before in \toyChunks {
    \draw[bracket] ({\boxleft + \start * \boxsize + 0.06},          {\boxbase - 0.08})
                -- ({\boxleft + \start * \boxsize + 0.06},          {\boxbase - 0.18})
                -- ({\boxleft + (\start + \len) * \boxsize - 0.06}, {\boxbase - 0.18})
                -- ({\boxleft + (\start + \len) * \boxsize - 0.06}, {\boxbase - 0.08});
  }%
}
